\section{Atmospheric Data Assimilation}

Data assimilation combines observations and numerical models
to produce an optimal estimate of the atmospheric state.

The general formulation is:

\begin{equation}
x_a = x_b + K(y-H(x_b))
\end{equation}

where:

\begin{itemize}
\item $x_a$ is the analysis state
\item $x_b$ is the background forecast
\item $y$ represents observations
\item $H$ is the observation operator
\item $K$ is the gain matrix
\end{itemize}



\section{Observation Systems}

ACF integrates multiple observation sources:

\begin{itemize}
\item Satellite measurements
\item Weather radar
\item Radiosondes
\item Surface meteorological stations
\item Aircraft observations
\item Ocean observations
\end{itemize}



\section{Variational Data Assimilation}


\subsection{3DVAR}

The 3DVAR method minimizes a cost function:

\begin{equation}
J(x)=
\frac{1}{2}(x-x_b)^TB^{-1}(x-x_b)
+
\frac{1}{2}(y-H(x))^TR^{-1}(y-H(x))
\end{equation}


where:

\begin{itemize}
\item $B$ is the background error covariance
\item $R$ is the observation error covariance
\end{itemize}



\subsection{4DVAR}

4DVAR introduces the time dimension:

\begin{equation}
J(x)=
J_b+
\sum_{i=0}^{n}
J_i
\end{equation}


It optimizes the atmospheric trajectory over a time window.



\section{Ensemble Methods}

Ensemble Kalman Filters estimate uncertainty using multiple
forecast realizations.

Applications:

\begin{itemize}
\item Probabilistic forecasting
\item Uncertainty estimation
\item Extreme event prediction
\end{itemize}
